% Article: The Path to Automated Narrative Processing
\documentclass[11pt,a4paper]{article}

% Packages
\usepackage[utf8]{inputenc}
\usepackage[T1]{fontenc}
\usepackage{geometry}
\usepackage{graphicx}
\usepackage{booktabs}
\usepackage{hyperref}
\usepackage{amsmath}
\usepackage{natbib}
\usepackage{setspace}

% Page geometry
\geometry{margin=1in}

% Line spacing - adjust for submission requirements
\onehalfspacing

% Hyperlinks
\hypersetup{
    colorlinks=true,
    linkcolor=blue,
    citecolor=blue,
    urlcolor=blue
}

% Metadata
\title{The Path to Automated Narrative Processing: \\
       Narrative and the Quest for Human-like Artificial Intelligence}
\author{Your Name\\
        \textit{Your Institution}\\
        \texttt{your.email@institution.edu}}
\date{\today}

\begin{document}

\maketitle

\begin{abstract}
How does narrative figure into the drive towards artificial intelligences that bear comparison to human cognition? There is a robust literature about narrative computation, focusing on the analysis of narrative structures. Large language models can be leveraged to generate narratives. While these two directions automate two high-level human cognitive activities, the processes do not resemble how narrative structures are recognized and created in human thinking. This article describes the function of ongoing narrative processing in human intelligence, followed by consideration of the state of artificial intelligence in approximating that functionality. We consider the problem ongoing narrative processing helps us solve, examine human event processing and current AI approaches, and offer a roadmap of research directions toward human-like event processing in AI. Finally, we address questions of embodiment, processing limitations, and the ethics of this research direction.
\end{abstract}

\textbf{Keywords:} narrative processing, artificial intelligence, event cognition, AGI, narrative computation

%=============================================================================
\section{Introduction}
\label{sec:introduction}
%=============================================================================

How does narrative figure into the drive towards artificial intelligences that bear comparison to human cognition?

There is a robust literature about narrative computation, focusing on the analysis of narrative structures. Large language models can be leveraged to generate narratives. While these two directions automate two high-level human cognitive activities, the processes do not resemble how narrative structures are recognized and created in human thinking.

This article describes the function of ongoing narrative processing in human intelligence, followed by consideration of the state of artificial intelligence in approximating that functionality.

%=============================================================================
\section{The Problem Narrative Processing Solves}
\label{sec:problem}
%=============================================================================

Ongoing narrative processing helps us solve a fundamental cognitive problem: orienting ourselves towards, framing responses to, and retroactively parsing experience.

% Expand this section with detailed discussion

\subsection{Orienting Towards Experience}

% Content to be developed

\subsection{Framing Responses}

% Content to be developed

\subsection{Retroactive Parsing}

% Content to be developed

%=============================================================================
\section{Human Event Processing}
\label{sec:human-processing}
%=============================================================================

% Discuss event segmentation, event cognition research, etc.

%=============================================================================
\section{AI Event Processing: State of the Art}
\label{sec:ai-processing}
%=============================================================================

% Survey current approaches in AI

\subsection{Narrative Analysis Systems}

% Content to be developed

\subsection{Narrative Generation}

% Content to be developed

\subsection{Gaps and Limitations}

% Content to be developed

%=============================================================================
\section{A Research Roadmap}
\label{sec:roadmap}
%=============================================================================

% Outline research directions that might lead to human-like event processing

\subsection{Direction 1}

% Content to be developed

\subsection{Direction 2}

% Content to be developed

\subsection{Contributions to AGI}

% Content to be developed

%=============================================================================
\section{Ethics and Open Questions}
\label{sec:ethics}
%=============================================================================

\subsection{Is Embodiment Essential?}

Is embodiment an essential precondition of narrative processing?

% Discussion to be developed

\subsection{Limitation or Feature?}

Is the ``functionality'' of narrative processing actually a limitation imposed by processing capacity, and might artificial systems ultimately transcend that limitation?

% Discussion to be developed

%=============================================================================
\section{Conclusion}
\label{sec:conclusion}
%=============================================================================

% Summary and concluding remarks

%=============================================================================
% References
%=============================================================================

\bibliographystyle{apalike}
\bibliography{references}

\end{document}
